\section{Related Work}
\label{sec:related}

\textbf{Sources of tail latency.} Dean and Barroso~\cite{dean2013tail}
showed that, in fan-out services, the latency tail of individual servers
sets the latency users see. Delimitrou and Kozyrakis~\cite{delimitrou2018amdahl} extended Amdahl's
law to tail latency with queueing models. They showed that a tail target,
unlike an average target, shifts the best design toward larger
cores, so what drives the tail directly affects how machines are
provisioned. Li et al.~\cite{li2014tales} measured where the
excess tail of null-RPC, Memcached and Nginx servers on Linux comes from:
background processes, FIFO-violating scheduling, the mapping of NIC
interrupts and threads to cores, and power saving. Each source was isolated
by changing the configuration and rerunning. Treadmill~\cite{zhang2016treadmill}
showed that load testers distort tail measurements (client-side queueing,
closed-loop arrivals, wrong aggregation) and used quantile regression over
many runs to attribute tail changes to hardware and software factors.
Both attribute the tail to a \emph{configuration}, offline and over many
runs, not to individual requests while the system runs.

\textbf{Waiting and scheduler analysis.} wPerf~\cite{zhou2018wperf} traces
waiting events between threads and builds wait-for graphs to find events
that cap throughput, with about 5\% recording overhead. Lozi et
al.~\cite{lozi2016wasted} found Linux load-balancing bugs that left cores
idle while threads waited in runqueues. These works target throughput and
scheduler correctness. They do not ask which OS events delayed a
particular slow request.

\textbf{Kernel tracing with eBPF.} Safe in-kernel instrumentation goes back
to BPF~\cite{mccanne1993bpf} and DTrace~\cite{cantrill2004dtrace}. eBPF adds
a verifier, maps and ring buffers~\cite{kernelbpf,ebpfdocs,vieira2020ebpf}.
The BCC tools~\cite{bcc,gregg2019bpf} report runqueue latency
(\texttt{runqlat}, \texttt{runqslower}), softirq time and off-CPU time as
histograms per CPU, process or cgroup. Netflix uses the same runqueue
signal per container to detect noisy neighbours~\cite{netflix2024noisy}.
Their output is aggregated over time or over a task, so it cannot be joined
with the latency of a single request. \texttt{runqlat} also folds wakeup
and preemption waits into one histogram. Recording every scheduler event instead,
as \texttt{perf sched} or ftrace do, is complete but expensive: our own
per-event streaming mode costs 63--104$\times$ more tracer CPU per request
than in-kernel accounting (\S\ref{sec:overhead}).

\textbf{Request-level tracing and mitigation.} Pivot
Tracing~\cite{mace2015pivot} and Hindsight~\cite{zhang2023hindsight}
follow requests across components. Hindsight keeps cheap trace data and
retrieves it only for edge cases such as slow requests. Their spans stop at
the application boundary, so scheduler and interrupt time inside a span
cannot be seen. Caladan~\cite{fried2020caladan} shows that fast
queueing-delay signals can drive core reallocation at microsecond scale,
but it replaces the Linux scheduler with its own runtime.

\textbf{Programmable scheduling.} A second line of work changes the
scheduler itself to cut tail latency. Shinjuku~\cite{kaffes2019shinjuku}
preempts requests every few microseconds so short requests no longer wait
behind long ones. ghOSt~\cite{humphries2021ghost} delegates Linux
scheduling policy to user-space agents. \texttt{sched\_ext}~\cite{schedext}
now lets such policies be written as BPF programs in the mainline kernel.
These systems make the policy flexible, but none of them says \emph{which}
OS mechanism delayed a given request. Without that, an operator cannot tell
whether a new policy, a different CPU placement or an interrupt-affinity
change is the right fix.

\textbf{Gaps.} We found no low-overhead mechanism that, on an unmodified
Linux kernel, tags each request with the scheduler and interrupt delay it
(and the requests queued ahead of it) suffered, compares that tag against
aggregate event counts and a queueing-only control, tests generalization to
an unseen workload, and validates the attribution with one-at-a-time
interventions. That is the gap this work addresses.
